\subsection{Inflationary QCD phase diagram --- arXiv:2512.24024}

\textbf{Authors:} Kohei Fujikura, Toshifumi Noumi

\paragraph{主な主張}
ローリングするインフラトンが誘起するフレーバー普遍な軸性化学ポテンシャルをもつ de Sitter 時空上の多フレーバー QCD を NJL 模型で解析し、十分大きな軸性化学ポテンシャルではインフレーション中またはその終了時に一次のカイラル相転移が起こり得ることを示した \cite{Fujikura:2025inflationary}。

\paragraph{新規性}
cosmological collider の対象を粒子スペクトルだけでなく物質の相構造へ拡張し、Hubble パラメータとインフラトン誘起軸性化学ポテンシャルで特徴付けられる ``inflationary QCD phase diagram'' を具体的なベンチマークとして構成した点にある。

\paragraph{位置付け}
インフレーションを高エネルギー粒子の探索装置として用いる cosmological collider と QCD 相転移研究を接続する研究であり、ローリングするスカラー場とフェルミオンの微分結合が初期宇宙の相構造そのものを変え得ることを示す具体例である。
